\chapter{Resultados}
\label{ch:resultados}

\section{Cobertura da base marítima}
\label{sec:inventario-resultados}

A matriz marítima contém 177 ligações com sentido definido entre dois portos. Em cada ligação, um porto é o ponto de partida e o outro é o ponto de chegada. Santos--Manaus e Manaus--Santos são ligações diferentes.

Para construir essas ligações, o sistema lê as tabelas de Carga e Atracação da ANTAQ. Ele organiza por data as escalas de cada viagem. Assim, identifica a ordem em que o navio visitou os portos. Os embarques e desembarques permitem reconstruir a carga a bordo entre duas escalas consecutivas.

Na base histórica, a tabela de Atracação forneceu 7.103 chamadas portuárias. Uma chamada é uma ocorrência de entrada de navio em porto. O sistema ordenou essas chamadas por navio e por data. Chamadas consecutivas do mesmo navio no mesmo porto foram reunidas em uma parada. O resultado contém 6.797 paradas distribuídas em 1.324 viagens reconstruídas.

Duas paradas consecutivas formam um subtrecho. Antes da padronização final dos portos, as 1.324 viagens continham 5.473 subtrechos candidatos, pois uma viagem com (n) paradas possui (n-1) subtrechos. Depois de reunir 35 transições que apontavam para o mesmo porto padronizado, restaram 5.438 subtrechos calculados.

Para analisar uma ligação, o sistema escolhe os dois portos, examina uma viagem de cada vez e lê suas paradas da mais antiga para a mais recente. O recorte começa quando o navio está no porto de partida e termina quando chega depois ao porto de destino, durante a mesma viagem. O navio pode ir diretamente ou parar em outros portos.

Na viagem real \codecell{voyage\_9612791\_00011}, o recorte Santos--Manaus começa na primeira parada e termina na quarta. Ele contém três subtrechos: Santos--Suape, Suape--Pecém e Pecém--Manaus. O nome técnico desse recorte é \emph{recorte histórico viagem--OD}, em que OD significa origem--destino. Cada viagem contribui uma única vez para a mesma ligação.

Uma viagem pode gerar vários recortes. Dentro de \codecell{voyage\_9612791\_00011}, o prefixo Santos--Suape--Pecém--Manaus gera seis combinações com sentido definido: Santos--Suape, Santos--Pecém, Santos--Manaus, Suape--Pecém, Suape--Manaus e Pecém--Manaus. A viagem completa ainda retorna de Manaus a Santos e pode fornecer outras ligações. Por isso, 1.324 viagens geraram 12.025 recortes históricos viagem--OD. A sequência completa de portos de cada recorte recebe o nome de \emph{corredor}. Em cada subtrecho, o sistema registra a distância, a carga histórica a bordo e a fonte do indicador de consumo.

O sistema faz quatro verificações antes de usar um recorte na estatística. Primeiro, a parte aproveitada da viagem deve seguir a direção pedida. No cálculo Santos--Manaus, ela começa na escala de Santos e termina na escala de Manaus; uma sequência Manaus--Suape--Santos não serve. Segundo, nenhum dos trechos navegados entre esses portos pode estar ausente. Terceiro, deve existir uma intensidade obtida pelo IMO ou estimada e identificada como fallback. Quarto, o trabalho de transporte deve ser maior que zero para o recorte receber peso na mediana; esse trabalho é a soma da carga a bordo multiplicada pela distância de cada trecho. Um recorte com resultado igual a zero permanece registrado, mas não altera a mediana quando existem resultados positivos. Se todos forem iguais a zero, o sistema usa a mediana simples e registra a exceção.

O indicador marítimo informa quantos gramas de combustível são consumidos para transportar uma tonelada por uma milha náutica. Sua unidade é g/(t$\cdot$nm), e o método o chama de \emph{intensidade}. Dos 5.438 subtrechos, 3.290 receberam uma intensidade específica porque o número IMO do navio foi encontrado no EU MRV.

Nos outros 2.148 subtrechos, o IMO não forneceu um valor específico. O sistema estimou a intensidade com embarcações semelhantes. Primeiro, procura o grupo mais específico disponível, chamado de classe. Se esse grupo não fornecer um valor, usa uma categoria mais ampla, chamada de tipo, como ``navio porta-contêineres''. Essa estimativa de grupo recebe o nome de \emph{fallback}. Na execução analisada, todos os 2.148 subtrechos substitutos foram resolvidos pelo tipo; nenhum recebeu fallback de classe. A regra de classe permanece implementada para bases que tragam esse metadado.

\begin{table}[htbp]
\centering
\small
\caption{Cobertura da matriz marítima baseada em ANTAQ e EU MRV.}
\label{tab:cobertura-matriz-maritima}
\begin{tabularx}{\textwidth}{L{0.42\textwidth}r Y}
\toprule
Indicador & Valor & Como ler o indicador \\
\midrule
Ligações entre portos com sentido definido & 177 & Relações para as quais existe um valor de intensidade marítima disponível.\\
Recortes históricos viagem--OD & 12.025 & Partes de uma viagem entre uma parada inicial e outra parada posterior.\\
Subtrechos navegados & 5.438 & Pernas consecutivas usadas para reconstruir carga, distância e consumo.\\
Subtrechos com intensidade por IMO & 3.290 & 60,50\% dos subtrechos; associação direta do navio ao EU MRV.\\
Subtrechos com fallback por tipo & 2.148 & 39,50\% dos subtrechos; estimativa robusta de grupo com fonte registrada.\\
\bottomrule
\end{tabularx}
\end{table}

Para calcular o fallback por tipo, o sistema retira 1\% dos menores valores e 1\% dos maiores quando a amostra permite. Depois, calcula a média do restante. Essa regra recebe o nome de média aparada. Em uma amostra pequena, na qual 1\% não retiraria nenhum valor, usa-se a mediana.

O fallback por classe usa a média aparada disponível no arquivo de classes. Se ela não estiver disponível, usa a mediana. Os valores obtidos diretamente pelo IMO não são retirados nem alterados. Em todos os casos, o resultado registra qual regra foi usada.

\section{Como obter um único indicador entre dois portos}
\label{sec:intensidade-od-resultados}

Esta etapa usa as viagens históricas da ANTAQ. Ela não usa ainda a massa informada pelo usuário. No cálculo Santos--Manaus, o sistema obtém uma intensidade para a parte de cada viagem que começa na escala de Santos e termina na escala de Manaus. Entram a ligação direta e as viagens que pararam em outros portos; Manaus--Suape--Santos não entra. Depois, reúne todos os recortes aprovados para obter um único valor para a ligação. Esse valor recebe o nome de \emph{intensidade representativa}.

Para saber quanto peso cada recorte histórico deve ter no cálculo, o sistema mede a atividade realizada. Primeiro, identifica a carga a bordo em cada subtrecho. Depois, multiplica essa carga pela distância navegada, em milhas náuticas. Por fim, soma os produtos de todos os subtrechos entre os dois portos. O resultado é expresso em t$\cdot$nm e recebe o nome de \emph{trabalho de transporte}.

O peso de cada recorte é a soma da carga a bordo multiplicada pela distância de cada trecho. Um recorte com maior soma recebe mais peso. O sistema ordena as intensidades da menor para a maior e acumula esses pesos. Escolhe a primeira intensidade em que a soma alcança 50\% do total. Essa regra recebe o nome de \emph{mediana ponderada pelo trabalho de transporte}.

Um recorte com trabalho igual a zero não altera o resultado quando existem pesos positivos. Se todos os recortes tiverem trabalho igual a zero, o sistema usa a mediana simples e registra essa exceção.

É importante separar duas perguntas. A primeira é: ``quais viagens históricas ajudam a calcular o indicador de consumo entre os dois portos?'' Para responder, o sistema examina uma viagem de cada vez. Ele cria o recorte entre a parada no porto de partida e uma parada posterior no porto de chegada. Não é necessário que a viagem tenha sido direta. O recorte também entra se houver uma ou várias paradas intermediárias.

A segunda pergunta é: ``qual corredor fornece a distância aplicada à carga informada pelo usuário?'' O sistema considera somente corredores inteiros observados dentro de uma única viagem, com distância conhecida em todos os subtrechos e um valor de consumo disponível. Se existir um recorte direto, usa esse corredor. Caso contrário, escolhe o corredor completo mais curto entre os que possuem escalas. O corredor escolhido recebe o nome técnico de \emph{caminho}. O sistema nunca une subtrechos de viagens diferentes para criar uma rota artificial.

\section{Estudo da ligação Santos--Manaus}
\label{sec:santos-manaus-resultados}

Santos--Manaus é um exemplo concreto dessa separação. A base histórica contém 89 recortes nos quais o navio saiu de Santos e chegou depois a Manaus na mesma viagem, sem faltar nenhum subtrecho intermediário e com os dados necessários para calcular a intensidade e o peso estatístico. Eles estão distribuídos em 22 corredores. Entre eles estão o recorte direto Santos--Manaus; Santos--Suape--Pecém--Manaus em \codecell{voyage\_9612791\_00011}; Santos--Rio de Janeiro--Suape--Pecém--Manaus em \codecell{voyage\_9784661\_00001}; e Santos--Navegantes--Pecém--Manaus em \codecell{voyage\_9852365\_00011}. O último não passa por Suape.

O corredor Santos--Suape--Pecém--Manaus aparece, por exemplo, em \codecell{voyage\_9612791\_00011}. Nessa viagem, as cargas a bordo foram 12.858,754~t, 16.718,333~t e 12.325,900~t nos três subtrechos. A intensidade veio diretamente do IMO 9612791 no EU MRV de 2023. A soma dos trabalhos de transporte foi \(39.294.668{,}494~\mathrm{t\,nm}\), associada a \(291.959{,}387\)~kg de combustível para a atividade histórica total do navio no recorte Santos--Manaus.

Todos esses corredores históricos alimentam a distribuição usada para calcular a intensidade. Entretanto, apenas o caminho selecionado fornece a distância aplicada à remessa do usuário. Como existe um recorte direto observado, o sistema escolhe Santos--Manaus. A matriz marítima atribui 6.112~km ou 3.300,216~nm ao subtrecho; essa distância é modelada, não uma trajetória medida do navio. Isso não exclui da estatística os recortes históricos com escalas.

\begin{table}[htbp]
\centering
\small
\caption{Indicadores de intensidade marítima do par Santos--Manaus.}
\label{tab:santos-manaus-intensidade}
\begin{tabularx}{\textwidth}{L{0.44\textwidth}r Y}
\toprule
Indicador & Valor & Como ler o indicador \\
\midrule
Recortes históricos usados no indicador & 89 & Partes de viagens que começam em Santos e terminam mais tarde em Manaus.\\
Corredores observados & 22 & Sequências completas e distintas de portos.\\
Caminho selecionado & Santos--Manaus direto & Sequência usada para obter a distância da remessa do usuário.\\
Distância selecionada & 6.112 km / 3.300,216 nm & Distância marítima aplicada à remessa.\\
Recortes com intensidade por IMO & 40 & Intensidade associada diretamente ao navio no EU MRV.\\
Recortes com fallback por tipo ou classe & 49 & Estimativa de grupo usada quando não houve correspondência por IMO.\\
Trabalho de transporte associado a fallback & 59,54\% & Participação dos fallbacks no peso estatístico total do par.\\
Mediana ponderada pelo trabalho de transporte & 9,322050 g/(t$\cdot$nm) & Intensidade representativa aplicada ao caminho selecionado.\\
Média aritmética ponderada diagnóstica & 25,243618 g/(t$\cdot$nm) & Indicador auxiliar; não é usado no cálculo do cenário.\\
\bottomrule
\end{tabularx}
\end{table}

A diferença entre 9,322050 e 25,243618 mostra que a distribuição tem intensidades elevadas capazes de puxar a média para cima. A mediana ponderada reduz essa influência sem apagar os recortes históricos. A média continua registrada como diagnóstico. A participação de 59,54\% do trabalho de transporte em fallback também permanece visível, pois limita a confiança que pode ser atribuída ao valor representativo.

O ponto de 50\% foi alcançado na viagem \codecell{voyage\_9697002\_00002}. Antes desse recorte, o peso acumulado era 49,168\% do total; depois de acrescentar seus \(34.357.307{,}013~\mathrm{t\,nm}\), passou a 50,257\%. Essa observação usou a intensidade \(9{,}322050~\mathrm{g/(t\,nm)}\), que se tornou o valor representativo. A fonte registrada foi a média aparada em 1\% do tipo-padrão documentado \emph{container ship}, usado porque não havia correspondência individual aplicável pelo IMO nem classe ou tipo específico informado.

\subsection*{Exemplo de execução para uma carga de 14 t}

As cargas reconstruídas nos 89 recortes históricos serviram para calcular os pesos estatísticos. Elas não são a carga desta execução. No teste de execução do modelo, foi usada uma entrada controlada de 14~t. A intensidade de 9,322050~g/(t$\cdot$nm) vem do conjunto histórico. A distância de 3.300,216~nm vem da matriz marítima para o caminho direto observado em \codecell{voyage\_9612789\_00004}. O consumo de navegação da entrada de teste é:

\begin{equation}
F_{mar} = \frac{9{,}322050 \times 14 \times 3.300{,}216}{1000}
= 430{,}71~\text{kg de combustível}.
\end{equation}

A multiplicação produz gramas de combustível; a divisão por 1.000 converte o resultado para quilogramas. Esse valor corresponde somente à navegação. Ele não é a emissão total em \cooe{}, o custo total da viagem ou o resultado porta a porta. Para obter esses totais, a execução ainda considera acessos rodoviários, operações portuárias e os fatores de emissão e custo definidos para o cenário.

\section{Comparação com uma referência externa}
\label{sec:resultados-benchmark}

O sistema compara seus resultados com uma planilha produzida fora do projeto. Essa referência externa recebe o nome de \emph{benchmark}. A planilha contém 36 células. Quinze não permitem comparação numérica: seis representam um porto comparado com ele mesmo e nove possuem valor rodoviário nulo ou não positivo. Restam 21 ligações origem--destino com valores positivos e utilizáveis.

Nos 21 pares, o benchmark e o \CabotageLens{} apontam a mesma direção: a alternativa com cabotagem apresenta menor emissão. Porém, os valores não têm a mesma magnitude. Por isso, o benchmark serve como verificação direcional, e não como calibração do modelo.

\begin{table}[htbp]
\centering
\small
\caption{Síntese da comparação com o benchmark externo.}
\label{tab:resultados-benchmark}
\begin{tabularx}{\textwidth}{L{0.34\textwidth}L{0.25\textwidth}Y}
\toprule
Métrica & Valor & Como ler o indicador \\
\midrule
Células da matriz & 36 & Total de posições lidas na matriz 6~x~6.\\
Pares OD positivos e suportados & 21 & Casos que permitem comparação direcional.\\
Células não comparáveis & 15 & Seis pares iguais e nove valores rodoviários nulos ou não positivos.\\
Alinhamento direcional & 21/21 & As duas fontes indicam menor emissão para a alternativa com cabotagem.\\
Classificação & \makecell[l]{\texttt{same\_direction\_}\\\texttt{large\_gap}} & Mesma direção, mas diferença relevante de magnitude.\\
Validação calibrada de magnitude & 0 & Nenhum caso demonstra reprodução exata do valor externo.\\
\bottomrule
\end{tabularx}
\end{table}

\section{Síntese dos resultados}
\label{sec:sintese-resultados}

Os resultados devem ser lidos em três níveis. A matriz informa quantas viagens, subtrechos e fontes sustentam o cálculo. O caso Santos--Manaus mostra que todas as sequências de viagem ajudam a calcular a intensidade. Em contraste, uma única sequência observada fornece os portos e a distância aplicados à carga. O benchmark verifica apenas se a direção geral é coerente com uma referência externa.

Esses números dependem da carga, dos portos, das distâncias e da cobertura MRV. Também dependem do que o cálculo inclui e exclui. Esse limite recebe o nome de \emph{fronteira}. As emissões incluem apenas a operação \ttw{} de \cooe{}. Os valores monetários incluem somente os custos modelados. Portanto, os resultados não são inventários \wtw{}/LCA, cotações de frete ou prova de disponibilidade comercial. O Capítulo~\ref{ch:discussao} explica como usar essa evidência sem extrapolar seu alcance.
